% -----------------------------------------------------------------------
% ScholAR: Page-Preserving Retrieval-Augmented Grounding
%          for Scientific Paper Comprehension
%
% AAAI-27 Anonymous Submission
% Formatted against the official AAAI-27 Author Kit
% (aaai2027.sty / aaai2027.bst, AuthorKit27/).
%
% HOW TO COMPILE:
%   pdflatex scholar_aaai27.tex
%   bibtex scholar_aaai27
%   pdflatex scholar_aaai27.tex
%   pdflatex scholar_aaai27.tex
% -----------------------------------------------------------------------

\documentclass[letterpaper]{article} % DO NOT CHANGE THIS
\usepackage[submission]{aaai2027}  % DO NOT CHANGE THIS
% The serif, sans-serif, and monospaced fonts are loaded automatically by
% aaai2027.sty (newtxtext, helvet, courier). DO NOT add \usepackage{times},
% \usepackage{helvet}, \usepackage{courier}, or any other font package.
\usepackage[hyphens]{url}  % DO NOT CHANGE THIS
\usepackage{graphicx} % DO NOT CHANGE THIS
\urlstyle{rm} % DO NOT CHANGE THIS
\def\UrlFont{\rm}  % DO NOT CHANGE THIS
\usepackage{natbib}  % DO NOT CHANGE THIS AND DO NOT ADD ANY OPTIONS TO IT
\usepackage{caption} % DO NOT CHANGE THIS AND DO NOT ADD ANY OPTIONS TO IT
\frenchspacing  % DO NOT CHANGE THIS

% ---- Math & algorithms (recommended by the author kit for papers with algorithms) ----
\usepackage{amsmath}
\usepackage{amssymb}
\usepackage{mathtools}
\usepackage{algorithm}
\usepackage{algorithmic}

% ---- Tables (recommended by the author kit) ----
\usepackage{booktabs}
\graphicspath{{figures/}}

\pdfinfo{
/TemplateVersion (2027.1)
}

\setcounter{secnumdepth}{2} % Section + subsection numbers enabled: the paper uses
                             % Section~\ref{...} cross-references to subsections
                             % throughout, which need depth 2 to number distinctly.

% ---- Title ----
\title{ScholAR: Page-Grounded Retrieval-Augmented Generation\\for Scientific Document Comprehension}

% Anonymous submission: sole author is "Anonymous Submission", affiliations empty.
\author{
    Anonymous Submission
}
\affiliations{
}

% -----------------------------------------------------------------------
\begin{document}
\maketitle

% -----------------------------------------------------------------------
\begin{abstract}
% -----------------------------------------------------------------------
Understanding and querying scientific papers demands precise, evidence-grounded
comprehension that general-purpose language models struggle to provide reliably.
Existing PDF-chat systems hallucinate citations, discard document structure, and
cannot trace generated claims back to verifiable source pages.
We present \textbf{ScholAR} (\textbf{S}mart \textbf{C}ompanion for \textbf{H}olistic \textbf{O}rganization, \textbf{L}iterature \textbf{A}nalysis \& \textbf{R}esearch), a retrieval-augmented generation framework designed
specifically for scientific document comprehension that runs \textbf{entirely on-device},
with no cloud API in the loop for either text or vision inference.
ScholAR enforces strict page-boundary-preserving text segmentation, ensuring that
every retrieved chunk maps exactly to a single PDF page.
An indirect citation grounding mechanism decouples LLM generation from
page-number prediction, replacing fragile in-context page references with
verifiable evidence identifiers that the frontend resolves to precise PDF highlights.
Beyond single-document text QA, ScholAR extends to \textbf{visual grounding}
(routing figure/table questions to a local, natively multimodal model) and
\textbf{multi-document citation-chasing} (resolving and ingesting a paper's
own references so questions can be answered against both the anchor paper and
its cited works).
We introduce a manually annotated retrieval benchmark spanning 14 query--chunk
pairs across 3 landmark NLP papers, evaluated using Recall@$K$ and
Mean Reciprocal Rank (MRR).
Our BM25-primary retrieval system achieves Recall@5 of 0.929 and MRR of 0.788,
outperforming keyword-overlap baselines by 10.1 percentage points in MRR.
On an 18-case visual-grounding pilot benchmark spanning four reasoning types,
retrieval correctly routes to the target figure or table in 18/18 cases.
On a 10-cluster multi-document locality benchmark, ScholAR retrieves the correct
cited paper within the top 5 results 80\% of the time (MRR 0.356), though
top-1 precision (10\%) remains a clear open challenge, consistent with prior
findings that even large foundation models struggle at cross-paper localization.
\end{abstract}

% -----------------------------------------------------------------------
\section{Introduction}
% -----------------------------------------------------------------------
% Target length: ~0.75 columns (approx. 350--450 words)
%
% Paragraph structure:
%   P1: Hook -- why reading research papers is hard; the trust problem with
%       general-purpose LLMs that summarize without grounding.
%   P2: Existing solutions and their limitations.
%       (Generic PDF chat tools, hallucinated citations, no structured retrieval)
%   P3: The specific problem this paper addresses.
%       (Page-preserving chunking + lexical grounding + indirect citation)
%   P4: Our contributions (bulleted list below).
%   P5: Paper outline.
%
% Contributions:
\begin{itemize}
  \item We propose \textbf{ScholAR}, an end-to-end, \textbf{fully local} RAG framework
        for scientific document comprehension that enforces strict page-boundary-preserving
        text segmentation and performs both text and vision inference on-device.
  \item We introduce an \textbf{indirect citation grounding} mechanism that decouples
        LLM generation from page number prediction, reducing hallucinated citations.
  \item We extend single-document grounding to \textbf{visual question answering} over
        figures and tables and to \textbf{multi-document citation-chasing}, resolving
        and ingesting a paper's own references so cross-paper questions can be answered.
  \item We present quantitative benchmarks for all four capabilities
        (retrieval, faithfulness, visual grounding, multi-document locality),
        each with manually annotated cases and Recall@$K$/MRR-style metrics.
  \item We conduct controlled \textbf{ablation experiments} comparing keyword overlap,
        BM25-only, and BM25-primary-with-reranking retrieval architectures, and
        compare two local models (\texttt{qwen3.5:9b}, \texttt{gemma4:12b}) on
        visual grounding to characterize model choice for a fully on-device system.
\end{itemize}


% -----------------------------------------------------------------------
\section{Related Work}
\label{sec:related-work}
% -----------------------------------------------------------------------

\subsection{Retrieval-Augmented Generation}
% Cover: Lewis et al. (2020) RAG paper, ColBERT, BEIR benchmark,
% standard RAG pipelines and their PDF-specific limitations.
The standard RAG framework was introduced by Lewis et al.\ \cite{lewis2020rag}.
BM25 probabilistic retrieval was formalised by Robertson and Zaragoza \cite{robertson2009bm25}.
The BEIR benchmark \cite{thakur2021beir} provides a heterogeneous zero-shot evaluation suite for retrieval systems.

\subsection{Scientific Literature RAG Systems}
\label{sec:related-sciRAG}
PaperQA \cite{lala2023paperqa} was an early high-accuracy RAG agent for scientific
QA and literature synthesis, but relies on proprietary cloud LLMs for both
retrieval agency and generation.
OpenScholar \cite{asai2024openscholar} substantially outperforms PaperQA and
GPT-4o on literature-synthesis quality using an open retrieval pipeline, though
generation still runs through large hosted or self-hosted models rather than
a single fully local system.
SciRAG \cite{ding2025scirag} further improves over both OpenScholar and PaperQA
via citation-graph-aware symbolic reasoning and outline-guided synthesis.
None of these systems targets a fully on-device deployment: all either call a
cloud API or assume access to large (multi-tens-of-billions-of-parameters)
self-hosted models. ScholAR instead targets consumer hardware, running both
text and vision inference through a single sub-15B-parameter local model,
trading some raw capability for zero data exposure and zero API dependency.

\subsection{Scientific Document Understanding}
% Cover: ScholarBERT, SciDuet, DOC2PPT, SlideTailor (AAAI-26),
% S2ORC, ACL Anthology, Grobid, Nougat for structured parsing.
GROBID \cite{lopez2009grobid} and Nougat \cite{blecher2023nougat} provide structured
parsing of academic PDFs into machine-readable formats.
SlideTailor \cite{zeng2026slidetailor} generates personalised slides from scientific papers
but relies on raw PDF text extraction, the same bottleneck ScholAR addresses.
The S2ORC corpus \cite{lo2020s2orc} provides a large-scale resource of parsed academic papers.

\subsection{Visual Document Retrieval}
ColPali \cite{faysse2024colpali} and its successors index each document page as
a visual embedding, enabling page retrieval directly from image content without
OCR or layout-aware text extraction. This is an architecturally different bet
from ScholAR's approach: rather than embedding whole pages as images, ScholAR
keeps BM25 text retrieval as the primary signal and routes only figure/table-targeted
queries to a vision-capable model applied to a cropped region around the
identified figure. This keeps the common text-QA path fast and fully lexical
while still enabling grounded visual question answering when a query names or
implies a specific figure or table.

\subsection{Multi-Modal and Multi-Document Scientific QA}
M3SciQA \cite{li2024m3sciqa} is the closest existing benchmark to ScholAR's
multi-document extension: it pairs each anchor paper with its cited references
and requires both multi-modal and cross-document reasoning, evaluating 18
foundation models and finding that all of them substantially underperform
human experts at this task. ScholAR's multi-document locality benchmark
(Section~\ref{sec:multidoc-bench}) is smaller in scope but directly comparable
in spirit, and our results echo M3SciQA's central finding: cross-paper
localization is hard even when the correct reference is already available in
the retrieval pool.

\subsection{Question Answering over Documents}
% Cover: SQuAD, Natural Questions, TriviaQA, SciFact (fact verification),
% QuALITY (long document QA), QASPER (scientific paper QA).
QASPER \cite{dasigi2021qasper} is the closest existing benchmark to our task,
covering information-seeking QA over scientific papers.
SciFact \cite{wadden2020scifact} evaluates fact verification over biomedical literature.

\subsection{Citation Grounding and Faithfulness}
Gao et al.\ \cite{gao2023attributing} study attribution in LLM generation;
their analysis shows that citation quality degrades sharply without retrieval grounding.
Min et al.\ \cite{min2023factscore} propose FActScore, which decomposes generated text
into atomic claims and verifies each claim against a knowledge source, adopting
this decomposition strategy in our benchmark construction.
Laban et al.\ \cite{laban2022summac} introduce SummaC for consistency detection
in summarization, with two variants: SummaC-Conv (fine-tuned NLI) and SummaC-ZS
(zero-shot via sentence embeddings).
We implement SummaC-ZS because it does not require a fine-tuned NLI model
(e.g., DeBERTa-MNLI), making it reproducible without GPU infrastructure or API access.
Zha et al.\ \cite{zha2023alignscore} show that embedding-based alignment metrics
substantially outperform ROUGE-family metrics for hallucination detection;
this motivates our choice of SummaC-ZS over lexical alternatives.


% -----------------------------------------------------------------------
\section{Problem Formulation}
% -----------------------------------------------------------------------
% Target length: ~0.25--0.50 columns
%
% Formally define the task:
%   Given a scientific PDF D, a user question q, and a corpus of
%   extracted chunks C = {c_1, ..., c_n}, the system must:
%   (1) Retrieve a ranked subset C_q ⊆ C of relevant chunks.
%   (2) Generate a grounded answer a using only evidence from C_q.
%   (3) Associate each claim in a with a verifiable citation (chunk_id, page).
%
% Introduce notation used throughout the paper.

\begin{equation}
  \hat{a} = \underset{a}{\arg\max}\; P(a \mid q,\, \mathcal{C}_q)
  \quad \text{where} \quad
  \mathcal{C}_q = \text{Retrieve}(q, \mathcal{C}, k)
\end{equation}


% -----------------------------------------------------------------------
\section{Method}
% -----------------------------------------------------------------------

\subsection{Document Processing Pipeline}
Given a paper (an arXiv identifier or an uploaded PDF), ScholAR downloads or
accepts the file, extracts per-page text via PyMuPDF, applies page-preserving
chunking (below) and figure/table extraction (Section~\ref{sec:method-visual}),
and persists the results as \texttt{metadata.json}, \texttt{pages.json},
\texttt{chunks.json}, and \texttt{figures.json}. Figure~\ref{fig:pipeline}
summarises the flow from raw PDF to a grounded, citable answer.

\begin{figure}[t]
  \centering
  % \includegraphics[width=\columnwidth]{pipeline}
  \caption{ScholAR document processing and retrieval pipeline.
           The system extracts page-level text, applies page-preserving chunking,
           and serves grounded answers with indirect citation IDs.}
  \label{fig:pipeline}
\end{figure}

\subsection{Page-Preserving Chunking}
Standard character- or token-window chunking routinely splits a chunk across
a page boundary, breaking the one-to-one mapping between a retrieved chunk
and the PDF page a citation should highlight. ScholAR instead chunks
strictly within each page: Algorithm~\ref{alg:chunking} slides a window of
$W{=}1400$ target words with $\delta{=}120$-word overlap over each page's
token stream independently, so a chunk never spans two pages. A lightweight
regex pass additionally tags each chunk with a detected section title
(Abstract, Method, Results, \ldots) when a heading-like line is found near
the top of the page; the retriever uses this as a section-hint signal below.

\begin{algorithm}[t]
  \caption{Page-Preserving Chunking}
  \label{alg:chunking}
  \begin{algorithmic}[1]
    \REQUIRE Pages $P = \{p_1, \ldots, p_n\}$, target words $W$, overlap $\delta$
    \ENSURE Chunks $\mathcal{C}$
    \STATE $\mathcal{C} \leftarrow [\;]$
    \FOR{each page $p_i \in P$}
      \STATE $w \leftarrow \text{split}(p_i.\text{text})$
      \STATE $s \leftarrow 0$
      \WHILE{$s < |w|$}
        \STATE $e \leftarrow \min(s + W,\; |w|)$
        \STATE $\mathcal{C}.\text{append}(\text{Chunk}(w[s{:}e],\; p_i.\text{page}))$
        \STATE $s \leftarrow e - \delta$
      \ENDWHILE
    \ENDFOR
    \RETURN $\mathcal{C}$
  \end{algorithmic}
\end{algorithm}

\subsection{BM25-Primary Retrieval with Heuristic Reranking}
The production retriever ranks chunks with Okapi BM25~\cite{robertson2009bm25}
($k_1{=}1.4$, $b{=}0.72$) as the primary signal, then applies a small set of
additive heuristic boosts: a query-expansion dictionary maps research-specific
terms to related vocabulary (e.g., \emph{result} $\to$ \emph{accuracy, BLEU,
score}); an explicit page number in the query adds a page-hint boost; query
language implying a paper section (e.g., \emph{limitation}, \emph{results})
adds a section-hint boost when a candidate chunk's detected section matches;
and rhetorical phrase cues (\emph{we propose}, \emph{we introduce},
\emph{outperforms}) add a small boost, since claim-bearing sentences are more
likely to satisfy an information-seeking query. Visual-cue boosting for
figure/table chunks (Section~\ref{sec:method-visual}) uses the same additive
mechanism. We initially explored a design in which a dense/semantic score
dominated ranking; our own retrieval ablation (Table~\ref{tab:retrieval})
showed BM25-primary was the more reliable choice on this benchmark, so
semantic, section, and phrase signals are kept as light rerankers rather than
a primary signal, while dense retrieval is used only in the separate Hybrid
BM25+Dense+RRF configuration evaluated for faithfulness
(Section~\ref{sec:faith-bench}).

\subsection{Indirect Citation Grounding}
Free-form citation prompting --- asking the model to write, e.g., ``[p.\ 7]''
directly in its answer --- is fragile: the model has no reliable way to know
which page a claim actually came from and will readily hallucinate a
plausible-looking page number. ScholAR instead assigns each piece of
retrieved evidence a short-lived evidence identifier (\texttt{E1},
\texttt{E2}, \ldots) for the current turn, and the generation prompt
instructs the model to cite \emph{only} these evidence IDs, never a page
number directly. After generation, ScholAR discards any page number the
model wrote on its own (a direct violation of the instruction) and rewrites
each valid evidence-ID citation into a compact numbered reference
(\texttt{[1]}, \texttt{[2]}, \ldots) in citation order; the frontend resolves
each numbered reference back to its source chunk's page and highlights the
corresponding region of the rendered PDF. This decouples the model's
citation behavior from page-number prediction entirely: the model only has
to decide \emph{which} retrieved evidence supports a claim, while the
mapping from evidence ID to actual page and highlight region is deterministic
application logic rather than a model output.

\subsection{Visual Grounding}
\label{sec:method-visual}
When a query is judged to reference visual content (via cue terms such as
``figure,'' ``table,'' or ``plot,'' or an explicit figure/table number),
retrieval boosts figure and table chunks so they can compete against much
longer text chunks in the BM25 ranking. Each figure and table caption on a
page is detected via regex, and the region above the caption is rendered at
$3\times$ zoom and cached as a standalone image alongside the caption text.
A query naming an explicit figure or table number (e.g., ``explain figure 2'')
is matched directly against that figure's own label rather than relying solely
on lexical overlap: this matters because our tokenizer, like most BM25
implementations, only begins a token on a letter, so ``Figure 2'' and
``Figure 14'' both reduce to the token \texttt{figure} once digits are
dropped, making them otherwise indistinguishable to the ranker. When the
top-ranked chunk is a figure or table, ScholAR base64-encodes the cached image
and sends it, together with the caption and supporting text context, to the
same local Ollama model used for text generation (\texttt{qwen3.5:9b} and
\texttt{gemma4:12b} are both natively multimodal), rather than to a separate
vision-specific model or a cloud vision API. If the image is missing or the
model call fails, the system degrades to a caption-only extractive answer
rather than failing outright.

\subsection{Multi-Document Extension}
\label{sec:method-multidoc}
ScholAR resolves an anchor paper's own bibliography and lets the user ingest
individual cited works into the same study session. For arXiv papers, the
reference list is resolved via the Semantic Scholar API by arXiv ID; for
uploaded PDFs without a known arXiv identifier, references are resolved via
an S2 title search over the last several pages of the document, then
regex-matched against both numbered (\texttt{[N]}) and author--year
(\texttt{Vaswani et al., 2017}) in-text citation styles. Ingesting a reference
downloads its PDF (when available), applies the same page-preserving
chunking pipeline as the anchor paper, and tags every resulting chunk with a
\texttt{source\_paper\_id}. At query time, chunks from all currently-loaded
secondary papers are merged with the anchor paper's chunks into a single
retrieval pool, and each returned citation is labeled with its source paper so
answers can distinguish anchor-paper evidence from evidence drawn from a
cited work.

\subsection{Study Goal Generation}
ScholAR generates eight paper-specific study goals per document using a local
Ollama model (\texttt{qwen3.5:9b} or \texttt{gemma4:12b}).
Each goal has a title, a 2--4 sentence description grounded in paper text,
and 3--5 recursive subquestions whose evidence chunks are resolved by
a second BM25 retrieval pass.
When no LLM is available, a fully extractive fallback generates goals
directly from sentence patterns in the abstract and introduction.

\subsection{Faithfulness Evaluation: SummaC-ZS + Semantic Coverage}
\label{sec:cfs-method}
We evaluate citation faithfulness using the SummaC-ZS protocol~\cite{laban2022summac},
the zero-shot variant of SummaC that does not require a fine-tuned NLI model.
The key advantage of SummaC-ZS over ROUGE-family metrics is that it operates
at the \emph{semantic} level via dense embeddings, not at the lexical level,
making it robust to paraphrased hallucinations that ROUGE cannot detect.

\textbf{Tier 1: SummaC-ZS (Semantic Entailment).}
Each ground-truth claim $c$ is first decomposed into atomic sub-claims
$\{a_1,\ldots,a_m\}$ via sentence- and clause-level splitting
(FActScore-style~\cite{min2023factscore}).
For each atom $a_i$ and each retrieved chunk $\mathcal{C}_j$, we compute:
\begin{equation}
  s(a_i, \mathcal{C}_j) = \max_{t \in \text{sents}(\mathcal{C}_j)} \cos\!\bigl(\phi(a_i),\, \phi(t)\bigr)
  \label{eq:summac-zs}
\end{equation}
where $\phi(\cdot)$ is the \texttt{all-MiniLM-L6-v2} sentence encoder~\cite{reimers2019sbert}
loaded from the local HuggingFace cache (no network access required),
and the $\max$ is over all sentences in the chunk.
Atom $a_i$ is classified \textsc{Entailment} if
$\max_j s(a_i, \mathcal{C}_j) \geq \tau_e$ (threshold $\tau_e = 0.42$),
\textsc{Neutral} if $\geq \tau_n = 0.25$, and \textsc{Contradiction} otherwise.
Numeric claims containing values not present in any retrieved chunk are
down-graded from \textsc{Entailment} to \textsc{Neutral}.
\begin{equation}
  \mathrm{SummaC\text{-}ZS}(c, \mathcal{C}) =
  \frac{|\{a_i : \mathrm{label}(a_i) = E\}|}{m}
  \label{eq:nli-cfs}
\end{equation}

\textbf{Tier 2: SCR (Semantic Coverage Rate).}
Whole-claim vs.\ best-chunk cosine similarity (a coarser, complementary signal):
$\mathrm{SCR}(c, \mathcal{C}) = \max_j \cos(\phi(c), \phi(\mathcal{C}_j))$.

\textbf{Tier 3: KFP (Key Fact Presence).}
Fraction of numeric values and technical proper nouns from the claim
found verbatim in the top-5 retrieved chunks (exact lexical recall for
quantities and named entities, where semantics cannot substitute).

The combined score is:
\begin{equation}
  \mathrm{CFS} = 0.50 \times \mathrm{SummaC\text{-}ZS} + 0.30 \times \mathrm{SCR} + 0.20 \times \mathrm{KFP}
  \label{eq:cfs}
\end{equation}

\subsection{System Architecture}
ScholAR is built as a fully local system.
The backend is a FastAPI application that stores parsed PDFs as JSON files
on disk (\texttt{metadata.json}, \texttt{pages.json}, \texttt{chunks.json}).
The frontend is a Next.js single-page application with an embedded
pdf.js viewer that highlights cited pages in real time.
All inference, both text generation and figure/table visual question
answering, is performed locally via Ollama using a single natively
multimodal model (\texttt{qwen3.5:9b} or \texttt{gemma4:12b}); no cloud
API is used anywhere in the pipeline.
No user data ever leaves the local machine.
All experiments in this paper were run on a consumer laptop (Apple Silicon,
18\,GB unified memory) under Python 3.12, PyTorch, and Ollama, with
\texttt{all-MiniLM-L6-v2}~\cite{reimers2019sbert} loaded from a local
Hugging Face cache for the faithfulness pipeline; no GPU or cloud compute
was used for any reported result.


% -----------------------------------------------------------------------
\section{Evaluation}
% -----------------------------------------------------------------------

\subsection{Retrieval Benchmark}
We evaluate retrieval quality on a manually annotated benchmark of 14 cases
spanning three landmark NLP papers:
\textit{Attention Is All You Need} (1706.03762~\cite{vaswani2017attention}),
\textit{Retrieval-Augmented Generation} (2005.11401~\cite{lewis2020rag}),
and \textit{LLaMA} (2302.13971~\cite{touvron2023llama}).
Each case specifies a user-style question, a set of expected relevant chunk IDs,
and the expected source pages.
Query types include main idea, architecture, training details, result tables,
human evaluation, safety analysis, and page-hint queries.
We report Recall@$K$ ($K \in \{1, 3, 5\}$) and Mean Reciprocal Rank (MRR).

\subsection{Faithfulness Benchmark}
\label{sec:faith-bench}
For the faithfulness evaluation, we construct \textbf{51 oracle-claim cases} manually
annotated across three landmark NLP papers:
\textit{Attention Is All You Need}~\cite{vaswani2017attention} (15 cases),
\textit{RAG}~\cite{lewis2020rag} (12 cases), and
\textit{LLaMA}~\cite{touvron2023llama} (24 cases).
Each case contains a natural-language query, a ground-truth claim,
a manually verified supporting chunk ID, and a claim type label.
Claim types include: \texttt{result\_number} (13), \texttt{training\_detail} (8),
\texttt{technical\_claim} (11), \texttt{architecture\_detail} (10),
\texttt{formula} (1), \texttt{conceptual\_claim} (5),
\texttt{human\_eval} (2), and \texttt{environmental\_claim} (1).
All annotations were produced manually by the authors.

\subsection{Visual Grounding Benchmark}
\label{sec:visual-bench}
We construct an 18-case pilot benchmark spanning four reasoning types:
\texttt{visual\_understanding} (6), \texttt{data\_extraction} (5),
\texttt{comparison} (4), and \texttt{location} (3). Each case specifies a
figure- or table-targeted question and the expected figure/table label(s)
that should be retrieved. We report two signals: \emph{retrieval grounding}
--- whether the correct figure or table is ranked first, isolating the
visual-cue retrieval and figure-number matching mechanism
(Section~\ref{sec:method-visual}) from the vision model's reading accuracy
--- and a lightweight \emph{answer-quality} proxy (the fraction of non-trivial
content words from the question itself that reappear in the generated
answer), which does depend on the vision model. Under \texttt{gemma4:12b}, retrieval grounding is 18/18
and every case is answered directly by the vision model with mean answer
score 0.731 (no case fell back to the caption-only degraded path).
Table~\ref{tab:visual} reports retrieval results by reasoning type. We
deliberately keep this benchmark small and frame it as a pilot study: it
validates that the routing mechanism works and gives a first read on answer
quality, not that visual question answering is solved.

\begin{table}[t]
  \centering
  \caption{Visual grounding retrieval results (18 cases). All cases correctly
           route to the target figure or table.}
  \label{tab:visual}
  \setlength{\tabcolsep}{4pt}
  \begin{tabular}{lcc}
    \toprule
    \textbf{Reasoning Type} & \textbf{$N$} & \textbf{Retrieval R@5} \\
    \midrule
    Visual understanding & 6 & 1.000 \\
    Data extraction      & 5 & 1.000 \\
    Comparison           & 4 & 1.000 \\
    Location             & 3 & 1.000 \\
    \midrule
    \textbf{Overall}      & 18 & \textbf{1.000} \\
    \bottomrule
  \end{tabular}
\end{table}

\subsection{Multi-Document Locality Benchmark}
\label{sec:multidoc-bench}
We construct an 18-case multi-document benchmark in the style of
M3SciQA~\cite{li2024m3sciqa}: 15 \emph{locality} cases, each requiring the
system to identify which cited reference (not the anchor paper) contains the
answer, and 3 \emph{detail} cases answerable from the anchor paper alone as a
sanity check. Of the 15 locality cases, 10 reference works with a resolvable
arXiv identifier (\texttt{locality\_arxiv}, our primary metric); the remaining
5 reference pre-arXiv works or datasets (e.g., WMT, Dropout, Wikipedia) that
cannot currently be downloaded and always score zero. Table~\ref{tab:multidoc}
reports Recall@$K$ and MRR for each subset.

\begin{table}[t]
  \centering
  \caption{Multi-document locality results. \texttt{locality\_arxiv} (10 cases)
           is the primary metric; \texttt{detail} cases confirm anchor-paper
           retrieval is unaffected by the larger multi-paper pool.}
  \label{tab:multidoc}
  \setlength{\tabcolsep}{4pt}
  \begin{tabular}{lcccc}
    \toprule
    \textbf{Subset} & \textbf{$N$} & \textbf{R@1} & \textbf{R@5} & \textbf{MRR} \\
    \midrule
    All                & 18 & 0.222 & 0.611 & 0.364 \\
    Locality (all)     & 15 & 0.067 & 0.533 & 0.237 \\
    Locality (arXiv)   & 10 & 0.100 & \textbf{0.800} & 0.356 \\
    Detail (anchor)    & 3  & \textbf{1.000} & \textbf{1.000} & \textbf{1.000} \\
    \bottomrule
  \end{tabular}
\end{table}

The correct cited paper is usually recoverable within the top 5 results
(80\%) but rarely ranks first (10\%): with dozens of papers' chunks merged
into a single flat retrieval pool, a niche secondary paper must win on lexical
match alone against the anchor paper's typically larger, more topically
central chunk set. Section~\ref{sec:discussion-multidoc} discusses this
finding relative to M3SciQA's comparable observation that even large
foundation models struggle at this exact task.

To give this number honest context rather than leaving it unexplained, we
compute two reference bounds on the 10 \texttt{locality\_arxiv} cases
(average 8.0 candidate secondary papers per case): a \emph{random floor} ---
the expected recall of picking a candidate paper uniformly at random with no
retrieval signal at all --- and an \emph{oracle} upper bound, where retrieval
is restricted to only the anchor paper plus the single correct secondary
paper (i.e., the paper-selection problem is assumed solved, isolating
chunk-level retrieval precision \emph{within} the correct paper).
Table~\ref{tab:multidoc-bounds} reports all three.

\begin{table}[t]
  \centering
  \caption{Multi-document locality: ScholAR vs.\ a random-guess floor and an
           oracle upper bound (10 \texttt{locality\_arxiv} cases).}
  \label{tab:multidoc-bounds}
  \setlength{\tabcolsep}{4pt}
  \begin{tabular}{lccc}
    \toprule
    \textbf{System} & \textbf{R@1} & \textbf{R@5} & \textbf{MRR} \\
    \midrule
    Random floor (no signal)        & 0.125 & 0.625 & 0.340 \\
    ScholAR (flat retrieval)        & 0.100 & \textbf{0.800} & 0.356 \\
    Oracle (correct paper isolated) & \textbf{0.700} & \textbf{0.900} & \textbf{0.783} \\
    \bottomrule
  \end{tabular}
\end{table}

Two things stand out. First, ScholAR's actual R@1 (0.100) is \emph{below}
the random-guess floor (0.125) --- at the single-guess precision level, flat
retrieval across a large multi-paper pool is not just imperfect, it is
slightly worse than chance, almost certainly because the anchor paper's
larger, more BM25-friendly chunk set systematically wins the top rank
regardless of which secondary paper is actually most relevant. At R@5
(0.800 vs.\ a 0.625 floor), the same flat retrieval clearly does carry a real
signal above chance once enough slots are available for it to surface.
Second, the oracle bound (R@1 0.700, MRR 0.783) shows that most of the gap
between ScholAR and a perfect system is a \emph{paper-selection} problem, not
a chunk-level retrieval problem: once the correct paper is isolated,
within-paper retrieval is nearly as strong as the single-document benchmark
in Table~\ref{tab:retrieval}. This directly motivates the per-paper-aware
retrieval and RLM directions discussed next, rather than further tuning of
chunk-level ranking.

\subsection{Local Model Comparison}
\label{sec:model-comparison}
Retrieval and faithfulness scoring in ScholAR are model-independent (BM25 and
a fixed local sentence encoder, with no LLM call in either path), so the only
benchmark where the choice of local Ollama model is actually measurable is
visual grounding, where the vision-capable model both reads the figure and
writes the answer. We run the 18-case visual grounding benchmark
(Section~\ref{sec:visual-bench}) under two natively multimodal local models,
\texttt{gemma4:12b} and \texttt{qwen3.5:9b}, holding everything else fixed.
Table~\ref{tab:model-comparison} reports retrieval grounding and a
keyword-overlap answer-quality proxy for each. Both models retrieve the
correct figure in all 18 cases and never fall back to the caption-only
degraded path, but the smaller \texttt{qwen3.5:9b} (9B parameters) scores
higher on answer quality than the larger \texttt{gemma4:12b} (12B parameters)
on this benchmark. We do not read this as a general claim that smaller models
are better --- the benchmark is small and the scoring proxy is a lexical
heuristic, not an LLM-judge --- but it does show that going fully local does
not require the largest available model, and that model choice is worth
benchmarking empirically per task rather than assumed from parameter count
alone.

\begin{table}[t]
  \centering
  \caption{Visual grounding under two local models (18 cases each).
           Answer score is a keyword-overlap heuristic, not an LLM-judge score.}
  \label{tab:model-comparison}
  \setlength{\tabcolsep}{4pt}
  \begin{tabular}{lccc}
    \toprule
    \textbf{Model} & \textbf{Retrieval R@5} & \textbf{Answer Score} & \textbf{Fallback} \\
    \midrule
    \texttt{gemma4:12b} & 1.000 & 0.731          & 0/18 \\
    \texttt{qwen3.5:9b} & 1.000 & \textbf{0.812} & 0/18 \\
    \bottomrule
  \end{tabular}
\end{table}

\subsection{Does Reading the Image Help? A Caption-Only Ablation}
\label{sec:caption-ablation}
Rather than implement a full ColPali-style~\cite{faysse2024colpali} visual
embedding baseline, which would require new retrieval infrastructure this
benchmark is not built to compare against, we run a cheaper, more direct
ablation: send the identical prompt to the same local model \emph{without}
the figure image attached, and score the answer with the same heuristic.
Under \texttt{gemma4:12b}, caption-only scores \textbf{0.752} versus
\textbf{0.731} for full vision --- removing the image does not lower this
particular score.

We do not read this as evidence that visual reading is unnecessary. Many
benchmark questions ask about a value also restated in the paper's prose
near the figure, so a text-only answer can be correct without consulting the
image; and our answer-quality proxy rewards lexical overlap with the
\emph{question} rather than grounded correctness, so a model that reports
fine-grained visual details absent from the question wording (e.g., a
specific symbol connecting two boxes in a diagram, observed qualitatively in
Section~\ref{sec:method-visual}) is not credited for them. Properly
separating these effects requires benchmark cases whose answer is genuinely
absent from the paper's prose, plus a correctness-based judge rather than a
lexical-overlap proxy; we flag this as necessary follow-up work rather than
overstate what the current ablation supports.

\subsection{Baseline Systems}
For retrieval, we compare five configurations:
(1) \textbf{Keyword Overlap}: Jaccard-style term overlap with page-hint boost;
(2) \textbf{BM25-Only}: standard BM25 with page-hint boost;
(3) \textbf{BM25+Rerank (No PH)}: BM25 without page hints;
(4) \textbf{BM25+Rerank (PH)}: full system with page hints (default);
(5) \textbf{Dense-Only}: single-pass \texttt{all-MiniLM-L6-v2} cosine similarity,
no BM25, no rerank, no page hints --- the naive-RAG point in the
lexical-to-dense taxonomy~\cite{thakur2021beir}.
For faithfulness, we compare:
(6) \textbf{BM25-primary}: the ScholAR production retriever;
(7) \textbf{Hybrid BM25+Dense+RRF}: dense retrieval fused via Reciprocal Rank Fusion~\cite{cormack2009rrf}.

\subsection{Retrieval Results}
Table~\ref{tab:retrieval} shows the retrieval results.
BM25-primary outperforms keyword overlap on R@1, R@3, and MRR, and ties it on
R@5.
Page hints do not change the aggregate MRR on this benchmark but improve
retrieval for specific page-reference queries.
Dense-Only unexpectedly beats every lexical baseline (R@1 0.786, R@3/R@5
1.000, MRR 0.881), contrary to BEIR's finding that zero-shot dense
retrieval trails BM25~\cite{thakur2021beir}, likely an artifact of this
benchmark's small size and paraphrased queries. At faithfulness scale (51
cases, Table~\ref{tab:faithfulness}) hybrid, not dense-only, wins, so we do
not read this as grounds to change the production retriever.

\begin{table}[t]
  \centering
  \caption{Retrieval results on 14 manually annotated cases.
           PH = page hints. Best values in \textbf{bold}.}
  \label{tab:retrieval}
  \setlength{\tabcolsep}{4pt}
  \begin{tabular}{lcccc}
    \toprule
    \textbf{System} & \textbf{R@1} & \textbf{R@3} & \textbf{R@5} & \textbf{MRR} \\
    \midrule
    Keyword Overlap    & 0.571 & 0.786 & 0.929 & 0.687 \\
    BM25 Only          & 0.714 & 0.857 & 0.929 & 0.788 \\
    BM25+Rerank (No PH)& 0.714 & 0.857 & 0.929 & 0.788 \\
    BM25+Rerank (PH)   & 0.714 & 0.857 & 0.929 & 0.788 \\
    Dense-Only         & \textbf{0.786} & \textbf{1.000} & \textbf{1.000} & \textbf{0.881} \\
    \bottomrule
  \end{tabular}
\end{table}

\subsection{Faithfulness Results}
Table~\ref{tab:faithfulness} reports SummaC-ZS faithfulness results on the 51-case benchmark.
The hybrid BM25+Dense+RRF system achieves a Combined CFS of \textbf{0.829},
improving over BM25-primary (0.809) by 2.0 percentage points.
More significantly, hybrid retrieval improves SCHR@5 from 0.824 to \textbf{0.922},
a gain of 9.8 percentage points, confirming that dense semantic retrieval
recovers chunks that BM25 lexical matching misses.
BM25-primary classifies 48 of 51 cases (94.1\%) and Hybrid \textbf{49 of 51}
(96.1\%) as FAITHFUL under the $\geq 0.55$ strict entailment threshold.
Both systems' SummaC-ZS Tier~1 entailment (0.941 / 0.961) and KFP (0.966 / 0.979)
scores are high, with SCR (whole-claim cosine) remaining the primary
differentiating factor between systems.

\begin{table}[t]
  \centering
  \caption{SummaC-ZS faithfulness results on 51 oracle-claim cases.
           T1 = SummaC-ZS entailment fraction; SCR = whole-claim cosine similarity;
           KFP = Key Fact Presence; SCHR@$K$ = Supporting Chunk Hit Rate.}
  \label{tab:faithfulness}
  \setlength{\tabcolsep}{4pt}
  \begin{tabular}{lcc}
    \toprule
    \textbf{Metric} & \textbf{BM25-primary} & \textbf{Hybrid+RRF} \\
    \midrule
    SummaC-ZS / T1 (entailment)  & 0.941        & \textbf{0.961} \\
    SCR / T2 (cosine similarity) & 0.484        & \textbf{0.509} \\
    KFP / T3 (key fact recall)   & 0.966        & \textbf{0.979} \\
    \midrule
    \textbf{Combined CFS}        & 0.809        & \textbf{0.829} \\
    SCHR@1                       & 0.373        & \textbf{0.392} \\
    SCHR@3                       & \textbf{0.725} & 0.686 \\
    SCHR@5                       & 0.824        & \textbf{0.922} \\
    Faithful ($\geq$ 0.55 CFS)   & 48 / 51      & \textbf{49 / 51} \\
    \bottomrule
  \end{tabular}
\end{table}

\subsection{Ablation: CFS vs.\ Retrieval Depth}
Table~\ref{tab:rank-ablation} shows how Combined CFS grows monotonically as more
retrieved chunks are included. This confirms that evidence for complex scientific
claims is distributed across multiple pages and that returning Top-5 chunks is
necessary to achieve full faithfulness coverage.
The hybrid system leads at most depths, with the largest gap at Top-1
(BM25 0.658 vs.\ Hybrid 0.669); the two are within 0.3 points of each other at
Top-3, where BM25 briefly edges ahead, before hybrid pulls back into the lead
at Top-4 and Top-5.

\begin{table}[t]
  \centering
  \caption{Combined CFS at each retrieval depth ($N$ = 51 cases).
           Both improve monotonically; hybrid leads at every depth except
           a near-tie at Top-3.}
  \label{tab:rank-ablation}
  \setlength{\tabcolsep}{5pt}
  \begin{tabular}{cccc}
    \toprule
    \textbf{Chunks} & \textbf{$N$} & \textbf{BM25} & \textbf{Hybrid} \\
    \midrule
    Top-1 & 51 & 0.658 & \textbf{0.669} \\
    Top-2 & 51 & 0.757 & \textbf{0.776} \\
    Top-3 & 51 & \textbf{0.799} & 0.796 \\
    Top-4 & 51 & 0.791 & \textbf{0.817} \\
    Top-5 & 51 & 0.809 & \textbf{0.829} \\
    \bottomrule
  \end{tabular}
\end{table}

\subsection{Ablation: CFS by Claim Type}
Table~\ref{tab:type-ablation} breaks down Combined CFS by claim type under BM25-primary.
Environmental and human evaluation claims score highest (0.918, 0.882), as their
supporting text is verbally rich and semantically dense.
Training-detail claims (0.668) are penalised by multi-hop evidence spread across
sections; the two chronic BM25 failure cases (\texttt{attn\_training\_steps}:
0.242 UNFAITHFUL; \texttt{llama\_training\_tokens}: 0.356 PARTIAL) both involve
numerical claims spread across multiple non-contiguous page ranges.
A third BM25 failure, \texttt{rag\_trivia\_score} (0.235 UNFAITHFUL, a
\texttt{result\_number} claim), is retrieval-induced rather than a distributed-evidence
case: BM25 misses the supporting chunk entirely in its Top-5 for this query,
while Hybrid still recovers it (0.802 FAITHFUL) via the dense component ---
see the discussion of the retrieval-precision/compound-term-matching tradeoff
below.

\begin{table}[t]
  \centering
  \caption{CFS breakdown by claim type, BM25-primary system, $N$=51.}
  \label{tab:type-ablation}
  \setlength{\tabcolsep}{3pt}
  \begin{tabular}{lccc}
    \toprule
    \textbf{Claim Type} & \textbf{$N$} & \textbf{BM25 CFS} & \textbf{Hybrid CFS} \\
    \midrule
    Environmental claim  & 1  & \textbf{0.918} & \textbf{0.918} \\
    Human evaluation     & 2  & 0.882          & 0.882 \\
    Result number        & 13 & 0.811          & \textbf{0.860} \\
    Conceptual claim     & 5  & 0.847          & 0.854 \\
    Technical claim      & 11 & 0.846          & 0.856 \\
    Architecture detail  & 10 & 0.830          & 0.845 \\
    Formula              & 1  & 0.876          & 0.876 \\
    Training detail      & 8  & 0.668          & 0.674 \\
    \bottomrule
  \end{tabular}
\end{table}


% -----------------------------------------------------------------------
\section{Discussion}
% -----------------------------------------------------------------------

\textbf{Hybrid retrieval improves semantic recall.}
BM25-primary achieves Recall@5 of 0.929 and MRR of 0.788,
confirming that lexical matching over page-preserving chunks is a reliable
foundation for scientific document QA.
Our SummaC-ZS evaluation further reveals that hybrid dense retrieval adds
meaningful value: SCHR@5 improves from 0.824 to 0.922 (+9.8 pp)
and Combined CFS rises from 0.809 to 0.829 (+2.0 pp).
These gains are driven by the dense component surfacing semantically relevant
chunks even when exact BM25 query terms are absent.

\textbf{SummaC-ZS detects paraphrased entailment that ROUGE cannot.}
Our Tier~1 SummaC-ZS score (0.961 for both systems) is significantly higher
than a purely lexical approach would yield for the same cases.
The key distinction is that SummaC-ZS uses sentence-level cosine similarity
(via \texttt{all-MiniLM-L6-v2}) which captures semantic equivalence across
paraphrased formulations, a property that ROUGE-family metrics fail to exhibit
\cite{zha2023alignscore}.
The SCR tier (whole-claim cosine, 0.484 BM25 vs.\ 0.509 Hybrid) is the primary
differentiating signal between the two retrieval systems.

\textbf{Training-detail claims are the main failure mode.}
The two non-FAITHFUL cases (\texttt{attn\_training\_steps}, \texttt{llama\_training\_tokens})
are both training-detail claims where numerical evidence is distributed across
non-contiguous page ranges. This validates structured, math-aware extraction
as the highest-priority future work for the single-document pipeline.

\textbf{Retrieval depth matters.}
The rank ablation shows monotonic CFS improvement from Top-1 (BM25: 0.658, Hybrid: 0.669)
to Top-5 (0.809 / 0.829), confirming the ScholAR default of $K{=}5$ is well-motivated.

\textbf{Tokenizer fixes trade off single-document precision against
compound-term matching; hybrid retrieval absorbs the cost better than
BM25 alone.}
We fixed a tokenizer bug uncovered during multi-document testing: token
boundaries only start on a letter, so a fused compound proper noun (e.g.\ a
figure or paper name written as one capitalized word) and its space-separated
form reduce to different token sets, breaking exact-name retrieval.
The fix is necessary for Sections~\ref{sec:method-visual}
and~\ref{sec:method-multidoc}, but it perturbs BM25's length normalisation
globally: it cost BM25-only Recall@5 seven points (1.00 $\to$ 0.929) and
introduced one new faithfulness retrieval miss (\texttt{rag\_trivia\_score}).
Hybrid BM25+Dense+RRF was materially more robust --- its numbers barely
moved, and it still recovers the case BM25 now misses --- showing hybrid
retrieval's real robustness to exactly this kind of perturbation.

\textbf{Cross-paper localization is the multi-document pipeline's weak point.}
\label{sec:discussion-multidoc}
The multi-document locality benchmark (Section~\ref{sec:multidoc-bench})
shows a wide gap between Recall@5 (0.80) and Recall@1 (0.10): the correct
secondary paper is usually retrieved but rarely ranks first. We attribute
this to the flat retrieval design --- all loaded papers' chunks share one
BM25 pool with no per-paper normalisation, so a niche cited paper competes
directly against the anchor paper's larger, more central chunk set. This
mirrors M3SciQA~\cite{li2024m3sciqa}, where even large foundation models
underperform human experts at the same cross-paper localization task. We
view per-paper-aware retrieval/reranking and the Recursive Language Model
pattern~\cite{zhang2025rlm} as the most promising directions for closing
this gap.

\textbf{Baseline scope.} Section~\ref{sec:related-sciRAG} explains why
OpenScholar, PaperQA2, and SciRAG are discussed qualitatively rather than
run head-to-head: none targets a single consumer laptop, so a fair
comparison needs a resource-matched setup (e.g., PaperQA2 on a local Ollama
backend), which we leave to future work rather than an untested,
crippled substitute or a cloud-API run that would contradict our design
constraint.


% -----------------------------------------------------------------------
\section{Future Work}
% -----------------------------------------------------------------------

The training-detail failure mode identified in Section~\ref{sec:faith-bench}
(numerical evidence spread across non-contiguous pages) motivates structured
extraction as the highest-priority next step:
we plan to integrate Nougat~\cite{blecher2023nougat} for LaTeX-preserving
extraction and evaluate it against ColPali~\cite{faysse2024colpali}-style
visual page retrieval as an alternative to our current caption-and-crop
figure grounding.
A math-aware tokeniser that indexes LaTeX tokens
(e.g., \texttt{\textbackslash sigma}, \texttt{\textbackslash frac}) alongside
natural language terms remains future work for the BM25 retrieval layer,
though rendered math display in the answer UI is already implemented via KaTeX.
We also plan to evaluate the Recursive Language Model
inference pattern~\cite{zhang2025rlm} as an alternative to flat retrieval
for multi-document citation-chasing, where our current locality benchmark
shows the clearest room for improvement.

% -----------------------------------------------------------------------
\section{Conclusion}
% -----------------------------------------------------------------------

We presented ScholAR, a fully local retrieval-augmented system for grounded
scientific document comprehension, extending single-document QA to visual
grounding and multi-document citation-chasing via page-preserving
segmentation and indirect citation grounding. Across four benchmarks:
BM25-primary retrieval reaches Recall@5 of 0.929 (14 cases); our NLI-CFS
faithfulness pipeline shows hybrid BM25+Dense+RRF raising Combined CFS to
0.829 and SCHR@5 to 0.922 (51 cases), with training-detail claims hardest;
visual grounding routes correctly in all 18 pilot cases; and multi-document
locality reaches 80\% Recall@5 but only 10\% at top-1 --- a flat-retrieval
limitation shared with comparable systems, and the clearest open problem
for future work.


% -----------------------------------------------------------------------
% Acknowledgements (omit for blind submission)
% -----------------------------------------------------------------------
% \section*{Acknowledgements}


% -----------------------------------------------------------------------
% References — aaai2027.sty sets the bibliography style automatically;
% do not add a \bibliographystyle command.
% -----------------------------------------------------------------------
\bibliography{scholar_references}

\end{document}
